\chapter{Mạng Xã Hội Và Bẫy So Sánh}
\markboth{Mạng Xã Hội Và Bẫy So Sánh}{Social Media and the Comparison Trap}

\section{Mạng Xã Hội Là Sân Khấu, Không Phải Toàn Bộ Đời Sống}
\begin{truyen}{Mạng Xã Hội Là Sân Khấu, Không Phải Toàn Bộ Đời Sống}{Social Media Is a Stage, Not a Whole Life}
\chuhoa{L}{ên mạng, em thấy người ta luôn đẹp hơn, giỏi hơn, vui hơn, thành công hơn.}
\textit{(Online, everyone seems prettier, smarter, happier, and more successful.)}

Điều đó không lạ.
\textit{(That is not strange.)}

Sân khấu nào chả được chọn góc đẹp nhất trước khi bật đèn.
\textit{(Every stage is arranged for its best angle before the lights go on.)}

Một học sinh nhìn ảnh bạn bè đi chơi rồi thấy mình thua thiệt.
\textit{(A student sees friends traveling and feels deprived.)}

Nhưng em không thấy bức ảnh đó được chụp sau ba mươi lần chọn góc, không thấy hôm trước họ vừa cãi nhau, không thấy khoản nợ mà gia đình phải gánh.
\textit{(But she does not see the thirty attempts behind the photo, the argument the day before, or the debt the family carries.)}

Mạng xã hội không nói dối hoàn toàn.
\textit{(Social media is not a total lie.)}

Nó chỉ kể sự thật đã được cắt gọt.
\textit{(It is a heavily edited truth.)}

\end{truyen}

\begin{baihoc}
Dạy học sinh rằng hình ảnh đẹp không phải bằng chứng của cuộc đời đẹp. Khả năng phân biệt sân khấu với đời thật là một kỹ năng sống.
\textit{(Teach students that beautiful images are not proof of a beautiful life. Distinguishing stage from reality is a life skill.)}
\end{baihoc}

\begin{ghinhoanh}
\textbf{Short English to remember:}

Teach students that beautiful images are not proof of a beautiful life.

Distinguishing stage from reality is a life skill.

\end{ghinhoanh}

\begin{tuhocnhanh}
\textbf{Cau hoi tu hoc / Self-study questions}

1. Van de chinh la gi? \textit{What is the main problem?}

2. Cach xu ly tot hon la gi? \textit{What is a better action?}

3. Mot cau tieng Anh em muon nho la cau nao? \textit{Which English sentence do you want to remember?}
\end{tuhocnhanh}
\ngancach

\section{Nghiện Sự Công Nhận Là Chiếc Xiềng Mềm}
\begin{truyen}{Nghiện Sự Công Nhận Là Chiếc Xiềng Mềm}{Addiction to Validation Is a Soft Chain}
\chuhoa{C}{ó học sinh đăng bài xong cứ năm phút lại mở điện thoại xem có ai thả tim chưa.}
\textit{(Some students post something and then check every five minutes to see who liked it.)}

Mỗi lượt thích là một mẩu đường nhỏ.
\textit{(Each like is a tiny piece of sugar.)}

Ngọt, nhanh, và không nuôi nổi điều gì lâu dài.
\textit{(Sweet, fast, and unable to nourish anything lasting.)}

Khi giá trị bản thân bị buộc vào phản ứng của người khác, học sinh sẽ dần đánh mất quyền tự đánh giá mình.
\textit{(When self-worth is tied to others’ reactions, students slowly lose the ability to evaluate themselves.)}

Sự lệ thuộc này rất mềm.
\textit{(This dependence is soft.)}

Nó không ồn ào như thuốc hay cờ bạc, nhưng nó bẻ cong hành vi mỗi ngày: chụp gì, nói gì, mặc gì, sống thế nào.
\textit{(It is not loud like drugs or gambling, but it bends daily behavior: what to photograph, what to say, what to wear, how to live.)}

\end{truyen}

\begin{baihoc}
Học sinh cần được tập làm việc tốt ngay cả khi không ai nhìn. Tự trọng bắt đầu ở chỗ đó.
\textit{(Students need practice doing good work even when no one is watching. Self-respect begins there.)}
\end{baihoc}

\begin{ghinhoanh}
\textbf{Short English to remember:}

Students need practice doing good work even when no one is watching.

Self-respect begins there.

\end{ghinhoanh}

\begin{tuhocnhanh}
\textbf{Cau hoi tu hoc / Self-study questions}

1. Van de chinh la gi? \textit{What is the main problem?}

2. Cach xu ly tot hon la gi? \textit{What is a better action?}

3. Mot cau tieng Anh em muon nho la cau nao? \textit{Which English sentence do you want to remember?}
\end{tuhocnhanh}
\ngancach

\section{Tin Đồn Trên Mạng Cắn Rất Lâu}
\begin{truyen}{Tin Đồn Trên Mạng Cắn Rất Lâu}{Rumors Online Bite for a Long Time}
\chuhoa{N}{gày xưa một câu nói xấu chỉ loanh quanh trong lớp.}
\textit{(In the past, a nasty comment might circle inside one classroom.)}

Bây giờ một ảnh chụp màn hình có thể chạy khắp trường trong một tối.
\textit{(Now a screenshot can spread across the whole school in one evening.)}

Người tung tin thường chỉ nghĩ: gửi cho vui.
\textit{(The person sharing it often thinks: just for fun.)}

Nhưng người bị nhắc tên có thể mất ngủ hàng tuần.
\textit{(The person named in it may lose sleep for weeks.)}

Môi trường số làm cho sự tàn nhẫn trở nên dễ hơn vì người ta không phải nhìn thẳng gương mặt đau của người bị hại.
\textit{(Digital spaces make cruelty easier because people do not have to look directly at the hurt face of the person harmed.)}

Nhà trường dạy công thức, nhưng cũng phải dạy học sinh rằng một lần bấm gửi có thể để lại hậu quả lâu hơn cả một bài kiểm tra trượt.
\textit{(Schools teach formulas, but they must also teach students that pressing send once can cause damage longer than failing a test.)}

\end{truyen}

\begin{baihoc}
Đạo đức trên mạng không phải môn phụ. Với học sinh hôm nay, nó là kỹ năng an toàn cơ bản.
\textit{(Digital ethics is not an extra subject. For today’s students, it is a basic safety skill.)}
\end{baihoc}

\begin{ghinhoanh}
\textbf{Short English to remember:}

Digital ethics is not an extra subject.

For today’s students, it is a basic safety skill.

\end{ghinhoanh}

\begin{tuhocnhanh}
\textbf{Cau hoi tu hoc / Self-study questions}

1. Van de chinh la gi? \textit{What is the main problem?}

2. Cach xu ly tot hon la gi? \textit{What is a better action?}

3. Mot cau tieng Anh em muon nho la cau nao? \textit{Which English sentence do you want to remember?}
\end{tuhocnhanh}
\ngancach

\section{Cần Có Những Khoảng Không Màn Hình}
\begin{truyen}{Cần Có Những Khoảng Không Màn Hình}{You Need Spaces Without Screens}
\chuhoa{C}{ó em than mình không tập trung được nữa.}
\textit{(Some students complain that they can no longer focus.)}

Đọc vài dòng đã muốn cầm điện thoại.
\textit{(After a few lines, they want to grab the phone.)}

Não quen với nội dung ngắn, nhanh, giật mạnh sẽ thấy việc học sâu như một cực hình.
\textit{(A brain trained by short, fast, highly stimulating content will experience deep study as torture.)}

Không phải học sinh lười hơn.
\textit{(It is not always that students have become lazier.)}

Nhiều khi môi trường số đã bẻ não các em theo hướng khó ngồi yên.
\textit{(Often the digital environment has bent their minds toward restlessness.)}

Giải pháp không hào nhoáng: những khoảng không màn hình, những giờ học không điện thoại, những buổi đi bộ không tai nghe, những cuốn sách giấy đọc chậm.
\textit{(The solution is not glamorous: spaces without screens, study hours without phones, walks without earphones, paper books read slowly.)}

\end{truyen}

\begin{baihoc}
Muốn lấy lại khả năng tập trung, học sinh phải chấp nhận kỷ luật với thiết bị. Không có sự tập trung nào lớn lên giữa một cơn nghiện kích thích liên tục.
\textit{(To regain focus, students must accept discipline with devices. No real concentration grows inside constant stimulation addiction.)}
\end{baihoc}

\begin{ghinhoanh}
\textbf{Short English to remember:}

To regain focus, students must accept discipline with devices.

No real concentration grows inside constant stimulation addiction.

\end{ghinhoanh}

\begin{tuhocnhanh}
\textbf{Cau hoi tu hoc / Self-study questions}

1. Van de chinh la gi? \textit{What is the main problem?}

2. Cach xu ly tot hon la gi? \textit{What is a better action?}

3. Mot cau tieng Anh em muon nho la cau nao? \textit{Which English sentence do you want to remember?}
\end{tuhocnhanh}
\ngancach

\section{Chiếc Gương Đồng Và Bộ Lọc Ảnh}
\begin{truyen}{Chiếc Gương Đồng Và Bộ Lọc Ảnh}{The Bronze Mirror and the Photo Filter}
\chuhoa{N}{gày xưa người ta soi mình trong gương đồng, hình đã méo sẵn.}
\textit{(Long ago people looked into bronze mirrors, where the image was already distorted.)}

Bây giờ học sinh soi mình qua bộ lọc ảnh, hình còn méo hơn.
\textit{(Today students look through photo filters, where the image is distorted even more.)}

Khác thời đại nhưng giống bản chất: con người rất dễ lầm cái hình đã sửa với con người thật.
\textit{(Different eras, same essence: human beings easily mistake the altered image for the real person.)}

Một bạn nữ chụp hàng chục ảnh mới dám đăng một.
\textit{(One girl took dozens of photos before daring to post one.)}

Rồi sau đó nhìn mặt mộc của mình và thấy thất vọng.
\textit{(Afterwards she looked at her natural face and felt disappointed.)}

Khi công nghệ làm cho chuẩn đẹp trở nên giả tạo hơn, lòng tự ti cũng có thêm nhiên liệu.
\textit{(When technology makes standards of beauty more artificial, insecurity gains more fuel.)}

\end{truyen}

\begin{baihoc}
Mạng xã hội không chỉ chỉnh ảnh; nó còn chỉnh chuẩn trong đầu học sinh. Cần dạy các em nhận ra điều đó.
\textit{(Social media does not only edit photos; it edits standards inside students’ minds. They need to be taught to notice that.)}
\end{baihoc}

\begin{ghinhoanh}
\textbf{Short English to remember:}

Social media does not only edit photos; it edits standards inside students’ minds.

They need to be taught to notice that.

\end{ghinhoanh}

\begin{tuhocnhanh}
\textbf{Cau hoi tu hoc / Self-study questions}

1. Van de chinh la gi? \textit{What is the main problem?}

2. Cach xu ly tot hon la gi? \textit{What is a better action?}

3. Mot cau tieng Anh em muon nho la cau nao? \textit{Which English sentence do you want to remember?}
\end{tuhocnhanh}
\ngancach

\section{Đứa Trẻ Sống Bằng Thông Báo}
\begin{truyen}{Đứa Trẻ Sống Bằng Thông Báo}{The Child Living by Notifications}
\chuhoa{M}{ột học sinh kể rằng em thấy bồn chồn hễ điện thoại im lặng quá lâu.}
\textit{(A student admitted feeling restless whenever the phone stayed quiet too long.)}

Thông báo đến không chỉ mang tin.
\textit{(A notification brings not only information.)}

Nó mang cảm giác mình vẫn được nhớ tới.
\textit{(It also brings the feeling of still being remembered.)}

Khi không có gì hiện lên, em thấy mình như biến mất một chút.
\textit{(When nothing appears, the student feels a little as if disappearing.)}

Đó là lúc công cụ liên lạc đã chạm sang vùng định nghĩa giá trị bản thân.
\textit{(That is the moment when a communication tool has crossed into defining self-worth.)}

\end{truyen}

\begin{baihoc}
Nếu sự im lặng của điện thoại làm em hoảng, vấn đề không còn là thiết bị nữa. Nó đã thành câu chuyện về sự lệ thuộc cảm xúc.
\textit{(If the silence of a phone makes you panic, the issue is no longer the device. It has become emotional dependence.)}
\end{baihoc}

\begin{ghinhoanh}
\textbf{Short English to remember:}

If the silence of a phone makes you panic, the issue is no longer the device.

It has become emotional dependence.

\end{ghinhoanh}

\begin{tuhocnhanh}
\textbf{Cau hoi tu hoc / Self-study questions}

1. Van de chinh la gi? \textit{What is the main problem?}

2. Cach xu ly tot hon la gi? \textit{What is a better action?}

3. Mot cau tieng Anh em muon nho la cau nao? \textit{Which English sentence do you want to remember?}
\end{tuhocnhanh}
\ngancach

\section{Tin Vịt Nhanh Hơn Thầy Cô}
\begin{truyen}{Tin Vịt Nhanh Hơn Thầy Cô}{Rumors Travel Faster Than Teachers}
\chuhoa{M}{ột tin giả về giáo viên hay học sinh có thể chạy khắp trường trước cả giờ ra chơi kết thúc.}
\textit{(A false rumor about a teacher or student can spread across a school before recess is over.)}

Ngày xưa lời đồn đi bộ.
\textit{(In the past rumors walked.)}

Bây giờ nó đi bằng mạng, nên tốc độ phá hỏng danh dự cũng lớn hơn nhiều.
\textit{(Now they travel by networks, so the speed of reputational harm is far greater.)}

Một học sinh chuyển tiếp vì không chịu nổi áp lực từ tin đồn mà không ai chịu gỡ.
\textit{(One student transferred because the pressure of an uncorrected rumor became unbearable.)}

Kỹ năng kiểm chứng thông tin bây giờ không chỉ là kỹ năng học tập.
\textit{(The skill of verifying information is no longer only an academic skill.)}

Nó là kỹ năng đạo đức.
\textit{(It is an ethical one.)}

\end{truyen}

\begin{baihoc}
Trước khi bấm chia sẻ, học sinh cần học cách dừng lại. Một lần dừng đúng lúc có thể cứu người khác khỏi rất nhiều ngày tối tăm.
\textit{(Before pressing share, students need to learn how to pause. One timely pause can save someone many dark days.)}
\end{baihoc}

\begin{ghinhoanh}
\textbf{Short English to remember:}

Before pressing share, students need to learn how to pause.

One timely pause can save someone many dark days.

\end{ghinhoanh}

\begin{tuhocnhanh}
\textbf{Cau hoi tu hoc / Self-study questions}

1. Van de chinh la gi? \textit{What is the main problem?}

2. Cach xu ly tot hon la gi? \textit{What is a better action?}

3. Mot cau tieng Anh em muon nho la cau nao? \textit{Which English sentence do you want to remember?}
\end{tuhocnhanh}
\ngancach

\section{Bộ Não Mỏi Vì Cuộn Mãi}
\begin{truyen}{Bộ Não Mỏi Vì Cuộn Mãi}{The Brain Tired by Endless Scrolling}
\chuhoa{M}{ột học sinh bảo mình không hề làm gì nặng mà vẫn rất mệt.}
\textit{(A student said he had done nothing strenuous yet felt very tired.)}

Em chỉ nằm và lướt điện thoại hàng giờ.
\textit{(He had only lain down and scrolled for hours.)}

Lướt liên tục tưởng là nghỉ, nhưng thật ra não vẫn bị kéo đi bởi chuỗi kích thích ngắn không dứt.
\textit{(Continuous scrolling feels like rest, but the brain is still being dragged by an endless chain of short stimuli.)}

Kết quả là nghỉ xong mà không hồi được, học vào cũng không sâu nổi.
\textit{(The result is rest that does not restore, and study that cannot go deep.)}

Không phải mọi hình thức nghỉ đều phục hồi.
\textit{(Not every form of rest restores.)}

Có kiểu nghỉ chỉ làm mình đờ hơn.
\textit{(Some forms only make us duller.)}

\end{truyen}

\begin{baihoc}
Học sinh cần học phân biệt giữa giải trí và phục hồi. Hai thứ đó không phải lúc nào cũng trùng nhau.
\textit{(Students need to distinguish entertainment from recovery. The two are not always the same.)}
\end{baihoc}

\begin{ghinhoanh}
\textbf{Short English to remember:}

Students need to distinguish entertainment from recovery.

The two are not always the same.

\end{ghinhoanh}

\begin{tuhocnhanh}
\textbf{Cau hoi tu hoc / Self-study questions}

1. Van de chinh la gi? \textit{What is the main problem?}

2. Cach xu ly tot hon la gi? \textit{What is a better action?}

3. Mot cau tieng Anh em muon nho la cau nao? \textit{Which English sentence do you want to remember?}
\end{tuhocnhanh}
\ngancach

\section{Một Ngày Không Mạng}
\begin{truyen}{Một Ngày Không Mạng}{One Day Offline}
\chuhoa{M}{ột lớp được giao bài tập lạ: một ngày cuối tuần không mạng xã hội rồi viết lại cảm giác của mình.}
\textit{(One class received an unusual assignment: spend one weekend day without social media and then write about the experience.)}

Nhiều em tưởng sẽ rất dễ.
\textit{(Many thought it would be easy.)}

Đến trưa đã thấy tay ngứa ngáy vì muốn cầm máy.
\textit{(By noon their hands already twitched to pick up the phone.)}

Tối đến, một số em viết rằng mình thấy thời gian dài hơn, yên hơn, nhưng cũng trống hơn.
\textit{(By evening, some wrote that time felt longer, quieter, but also emptier.)}

Bài tập đó không nhằm tôn thờ đời sống không công nghệ.
\textit{(The exercise was not about glorifying a life without technology.)}

Nó chỉ giúp học sinh thấy mình đang lệ thuộc đến đâu.
\textit{(It simply helped students see how dependent they had become.)}

\end{truyen}

\begin{baihoc}
Muốn quản lý một thứ, trước hết phải thấy mức độ nó đang quản lý mình. Khoảng tách ra là cách nhìn rõ nhất.
\textit{(To manage something, you must first see how much it already manages you. Distance is often the clearest lens.)}
\end{baihoc}

\begin{ghinhoanh}
\textbf{Short English to remember:}

To manage something, you must first see how much it already manages you.

Distance is often the clearest lens.

\end{ghinhoanh}

\begin{tuhocnhanh}
\textbf{Cau hoi tu hoc / Self-study questions}

1. Van de chinh la gi? \textit{What is the main problem?}

2. Cach xu ly tot hon la gi? \textit{What is a better action?}

3. Mot cau tieng Anh em muon nho la cau nao? \textit{Which English sentence do you want to remember?}
\end{tuhocnhanh}
\ngancach

\section{Học Sâu Trong Thời Đại Hời Hợt}
\begin{truyen}{Học Sâu Trong Thời Đại Hời Hợt}{Deep Study in a Shallow Age}
\chuhoa{N}{gày xưa người ta sợ thiếu sách.}
\textit{(In the past people feared a lack of books.)}

Bây giờ học sinh chìm trong quá nhiều nội dung nhưng lại thiếu chiều sâu.
\textit{(Now students drown in too much content while lacking depth.)}

Biết rất nhiều mẩu vụn không giống với hiểu được một điều đến nơi đến chốn.
\textit{(Knowing many fragments is not the same as understanding one thing thoroughly.)}

Một thầy giáo bảo lớp: "Nếu em không còn đọc nổi mười trang liên tiếp, đó không chỉ là vấn đề của môn Văn.
\textit{(A teacher told the class, "If you can no longer read ten pages in a row, that is not only a literature problem.)}

Đó là vấn đề của sức bền trí tuệ."
\textit{(It is a problem of intellectual endurance.")}

Trong thời đại này, khả năng ngồi lâu với một vấn đề có thể trở thành lợi thế hiếm hơn cả điểm cao.
\textit{(In this age, the ability to stay with one problem for a long time may become rarer, and more valuable, than high scores.)}

\end{truyen}

\begin{baihoc}
Muốn học thật, học sinh phải bơi ngược lại thói quen tiêu thụ nhanh. Sự sâu sắc luôn tốn thời gian.
\textit{(To learn for real, students must swim against the habit of rapid consumption. Depth always costs time.)}
\end{baihoc}

\begin{ghinhoanh}
\textbf{Short English to remember:}

To learn for real, students must swim against the habit of rapid consumption.

Depth always costs time.

\end{ghinhoanh}

\begin{tuhocnhanh}
\textbf{Cau hoi tu hoc / Self-study questions}

1. Van de chinh la gi? \textit{What is the main problem?}

2. Cach xu ly tot hon la gi? \textit{What is a better action?}

3. Mot cau tieng Anh em muon nho la cau nao? \textit{Which English sentence do you want to remember?}
\end{tuhocnhanh}
